\providecommand{\P}{\mathbb{P}}
\providecommand{\C}{\mathbb{C}}
\providecommand{\Z}{\mathbb{Z}}
\providecommand{\F}{\mathbb{F}}
\providecommand{\cO}{\mathcal{O}}
\providecommand{\cH}{\mathcal{H}}
\providecommand{\fgl}{\mathfrak{gl}}
\providecommand{\fsl}{\mathfrak{sl}}
\providecommand{\CoHA}{\mathrm{CoHA}}
\providecommand{\Tot}{\mathrm{Tot}}
\providecommand{\Crit}{\mathrm{Crit}}
\providecommand{\tr}{\mathrm{tr}}
\providecommand{\Spec}{\mathrm{Spec}}
\providecommand{\Coh}{\mathrm{Coh}}
\providecommand{\End}{\mathrm{End}}
\providecommand{\Hom}{\mathrm{Hom}}
\providecommand{\Jac}{\mathrm{Jac}}
\providecommand{\Pic}{\mathrm{Pic}}
\providecommand{\Hilb}{\mathrm{Hilb}}

\section*{Local Hirzebruch $K_{\F_n}$ and the Hall positive half}
\label{sec:local-hirzebruch-yplus}

Fix the Hirzebruch surface
\[
 \F_n=\P(\cO_{\P^1}\oplus\cO_{\P^1}(-n)),\qquad n\geq0,
\]
with negative section \(s\), fibre \(f\), and intersection form
\[
 s^2=-n,\qquad s\cdot f=1,\qquad f^2=0.
\]
Its canonical bundle is
\[
 K_{\F_n}=-2s-(n+2)f,\qquad K_{\F_n}^2=8,
\]
and its total space
\[
 X_n=K_{\F_n}=\Tot(\omega_{\F_n})
\]
is a noncompact local Calabi--Yau threefold with compact core
\(\F_n\). The local chiral scalar belongs to this compact-core
normalisation:
\[
 \kappa_{\mathrm{ch}}^{\mathrm{loc}}(X_n)
 =\frac{1}{2}\chi_{\mathrm{top}}(\F_n)=2.
\]
This is a compact-core invariant. The compact Calabi--Yau
supertrace \(\kappa_{\mathrm{cat}}=\chi(\cO_X)\) has compactness
hypotheses outside the noncompact total space \(X_n\), since
\[
 \pi_*\cO_{X_n}=\bigoplus_{m\geq0}\omega_{\F_n}^{-m}
\]
has infinite rank over \(\F_n\). A Borcherds weight
\(\kappa_{\mathrm{BKM}}\) is also absent until a Lorentzian lattice,
Jacobi input, and denominator identity are supplied.

\subsection*{The toric fan}

The fan \(\Sigma_{\F_n}\subset\Z^2\) has rays
\[
 v_1=(1,0),\quad v_2=(0,1),\quad v_3=(-1,n),\quad v_4=(0,-1),
\]
with maximal cones \(\langle v_i,v_{i+1}\rangle\), indices taken
cyclically. The parameter \(n\) is the shear in \(v_3\). The cases
\(n=0\) and \(n=1\) are Fano:
\(\F_0=\P^1\times\P^1\) and
\(\F_1\simeq\mathrm{Bl}_{pt}\P^2\). The anticanonical class on
\(\F_2\) is nef with zero intersection \((-K_{\F_2})\cdot s=0\), hence
it lies on the boundary of the ample cone. For \(n\geq3\),
\(-K_{\F_n}\) fails nefness.

The fan of \(X_n\) lies in \(\Z^3\). Put
\[
 u_0=(0,0,1),\qquad \bar v_i=(v_i,1).
\]
The maximal cones are
\[
 \sigma_i=\langle u_0,\bar v_i,\bar v_{i+1}\rangle,\qquad i\in\Z/4\Z .
\]
Their determinants are the determinants of the smooth surface cones,
hence every affine chart is smooth. All primitive generators lie on
the affine hyperplane \(z=1\), the toric Calabi--Yau condition. The
divisor \(D_{u_0}\) has star \(\Sigma_{\F_n}\), hence
\[
 D_{u_0}\simeq\F_n
\]
is the compact zero section. Thus \(X_n\) has a compact four-cycle.
The Rapcak--Soibelman--Yang--Zhao compact-four-cycle-free theorem gives
the \(\C^3\) chart model here and leaves the global compact-divisor
sector as additional data.

\subsection*{Affine charts and transition cocycles}

For every \(i\),
\[
 U_i=\Spec\C[\sigma_i^\vee\cap M]\simeq\C^3,
\]
where \(M=\Hom(\Z^3,\Z)\). Adjacent charts meet along a smooth
two-dimensional cone, so
\[
 U_i\cap U_{i+1}\simeq\C^2\times\C^*.
\]
Opposite charts share only the ray \(u_0\), so
\[
 U_i\cap U_{i+2}\simeq\C\times(\C^*)^2.
\]
On the wall between \(v_2\) and \(v_3\), one base change contains the
integral shear
\[
 \begin{pmatrix} -1 & 0 & 0 \\ n & 1 & 0 \\ 0 & 0 & 1 \end{pmatrix}
 \in\mathrm{GL}_3(\Z),
\]
with determinant \(-1\). The logarithmic volume form is preserved up
to sign. The cocycle remembers \(n\), although every local chart is
\(\C^3\).

\subsection*{The surface collection and the \(3\)-Calabi--Yau completion}

Put \(H=s+nf\). A standard toric exceptional collection on \(\F_n\) is
\[
 E_\bullet(\F_n)=
 \bigl(\cO,\ \cO(f),\ \cO(H),\ \cO(H+f)\bigr).
\]
The elementary pushforward identity
\[
 \pi_*\cO_{\F_n}(s+bf)
 =\cO_{\P^1}(b)\oplus\cO_{\P^1}(b-n)
\]
gives the Hom dimensions
\begin{align*}
 &\dim\Hom(\cO,\cO(f))=2,\qquad
 \dim\Hom(\cO,\cO(H))=n+2,\qquad
 \dim\Hom(\cO,\cO(H+f))=n+4,\\
 &\dim\Hom(\cO(f),\cO(H))=n,\qquad
 \dim\Hom(\cO(f),\cO(H+f))=n+2,\qquad
 \dim\Hom(\cO(H),\cO(H+f))=2 .
\end{align*}
These are raw Hom dimensions. The arrows of a minimal quiver are the
indecomposable generators, i.e. Hom spaces modulo images of lower
compositions. The relations are the multiplication maps
\[
 H^0(\cO(D_1))\otimes H^0(\cO(D_2))
 \longrightarrow H^0(\cO(D_1+D_2)).
\]
Consequently the quiver is presentation data, while the dependence on
\(n\) is already forced by \(s^2=-n\) and by the summand
\(\cO_{\P^1}(b-n)\) in \(\pi_*\cO_{\F_n}(s+bf)\).

The local threefold category is obtained from the surface endomorphism
algebra by the \(3\)-Calabi--Yau completion. In strict quiver language
this is the Ginzburg completion of a quiver with potential
\((Q_n,W_n)\). In an \(A_\infty\)-presentation the same information is
carried by higher products and Serre duality. The opposite arrows,
loops, and cyclic derivatives are part of the completion. An arrow
count alone falls short of an invariant of \(X_n\).

\subsection*{Critical CoHA and the positive half}

For a chosen presentation, stability chamber \(\theta\), and
orientation datum, write
\[
 Y^+_{\mathrm{Hall}}(X_n;\theta)
 :=\CoHA(Q_n,W_n;\theta)
\]
for the critical cohomological Hall algebra defined with vanishing
cycles for \(\tr W_n\). It is an associative \(E_1\) Hall algebra.
The symbol \(Y^+\) means positive half.

The fixed-point calculation is a theorem. On every affine chart
\(U_i\simeq\C^3\), with the standard equivariant Calabi--Yau condition
on the torus weights, the critical CoHA of the Jordan three-loop model is
\[
 \CoHA_T(\C^3)=Y^+(\widehat{\mathfrak{gl}}_1).
\]
The global algebra requires more than a tensor product of four chart
algebras. The compact divisor \(D_{u_0}\simeq\F_n\), the transition
cocycle, the orientation data, and wall-crossing for \(\theta\) all
enter the global CoHA. A statement of the form
\[
 \CoHA(Q_n,W_n;\theta)\cong Y^+(\widehat{\mathfrak{g}}_{Q_n})
\]
requires an explicit \((Q_n,W_n)\), a stability chamber, and a PBW or
shuffle presentation. Before these data are proved, the honest object
is the chamberwise critical CoHA \(Y^+_{\mathrm{Hall}}(X_n;\theta)\).

The positive half precedes every doubled or centred object:
\[
 Y^+_{\mathrm{Hall}}(X_n;\theta)
 \quad\leadsto\quad
 D^\wedge\!\bigl(Y^+_{\mathrm{Hall}}(X_n;\theta)\bigr)
 \quad\leadsto\quad
 Z^{\mathrm{der}}_{\mathrm{ch}}
 \quad\leadsto\quad
 \text{possible \(W\)-algebra evaluation}.
\]
The CoHA itself is the positive Hall half. The Drinfeld double,
completion, derived chiral centre, and any full \(W\)-algebra
comparison are additional layers. Since \(X_n\) has Calabi--Yau
dimension \(3\), the \(\Phi_3\)-output in this note is \(E_1\)-chiral.
Braided \(E_2\) structure belongs to the centre, double, or module
category above the positive half.

\subsection*{Characters and the \(n\)-dependent refinement}

G\"ottsche's rank-one surface Hilbert-scheme formula gives
\[
 \sum_{m\geq0}\chi_{\mathrm{top}}\bigl(\Hilb^m(\F_n)\bigr)q^m
 = \prod_{r\geq1}(1-q^r)^{-4}.
\]
The exponent is \(\chi_{\mathrm{top}}(\F_n)=4\), obtained either from
the four torus fixed points or from \(b_0+b_2+b_4=1+2+1\). This
one-variable character is the source of the local-surface
\(\kappa_{\mathrm{ch}}^{\mathrm{loc}}\)-normalisation and is independent
of \(n\).

The refined Hall or topological-vertex character must keep the two
curve variables \(Q_s,Q_f\). There \(n\) enters through
\[
 s^2=-n,\qquad K_{\F_n}=-2s-(n+2)f,
\]
and through the chamber wall-crossing of \((Q_n,W_n)\). A
presentation-independent character of
\(Y^+_{\mathrm{Hall}}(X_n;\theta)\) must either retain these variables
or prove the reduction to the chosen one-variable specialisation.

\subsection*{\(n=0\): the product base}

For \(n=0\), \(\F_0=\P^1\times\P^1\) and
\[
 X_0=\Tot\bigl(\cO_{\P^1\times\P^1}(-2,-2)\bigr).
\]
The two rulings give two curve variables and a square toric diagram.
The total space is a line bundle over a product surface. A
tensor-product decomposition of \(Y^+_{\mathrm{Hall}}(X_0;\theta)\)
requires a K\"unneth theorem for critical CoHAs with potential and
compatible stability chambers. The fan gives the unconditional
statement: the two-ruling symmetry exchanges the \(s\)- and
\(f\)-directions, and
\[
 \kappa_{\mathrm{ch}}^{\mathrm{loc}}(X_0)=2.
\]

\subsection*{\(n=2\): the \(A_1\) local neighbourhood}

For \(n=2\), the negative section is a rational \((-2)\)-curve. The
open neighbourhood obtained by deleting the section at infinity is
\[
 \Tot(\cO_{\P^1}(-2))\longrightarrow \C^2/\Z_2,
\]
the minimal \(A_1\)-resolution. Since this surface is symplectic, its
canonical bundle is trivial. The restriction of \(X_2\) to this
neighbourhood is therefore the crepant local model over
\((\C^2/\Z_2)\times\C\), with stack form
\[
 [\C^2/\Z_2]\times\C .
\]
Bridgeland--King--Reid McKay applies to this local quotient chart. The
global threefold \(X_2=K_{\F_2}\) still contains the compact divisor
\(\F_2\), four toric fixed charts, and the local-surface wall-crossing
problem. The \(\Z_2\)-orbifold quiver is a local \(A_1\) contribution;
the global positive half requires the \(X_2\) potential and chamber
data.

\subsection*{The local chiral scalar and the \(n\)-dependent data}

The local-surface modular characteristic is
\[
 \kappa_{\mathrm{ch}}^{\mathrm{loc}}(X_n)
 :=\kappa_{\mathrm{ch}}(A_{\mathrm{loc}\,\F_n})
 =\frac{1}{2}\chi_{\mathrm{top}}(\F_n)
 =\frac{4}{2}
 =2 .
\]
This is the same local-surface normalisation that gives
\(\kappa_{\mathrm{ch}}(K_{\P^2})=3/2\), since
\(\chi_{\mathrm{top}}(\P^2)=3\). The resolved conifold is a different
noncompact CY$_3$: it is \(\Tot(\cO(-1)^{\oplus2}\to\P^1)\), a
curve-base geometry separate from local surfaces \(\Tot(K_S\to S)\),
and its separate value is
\(\kappa_{\mathrm{ch}}=1\) by direct McKay computation.

For \(X_n\), the number \(2\) is constant in \(n\). The parameter
\(n\) is detected by the intersection form, by the multiplication maps
in the surface presentation, and by the refined two-variable Hall
character. Any isomorphism
\[
 Y^+_{\mathrm{Hall}}(X_n;\theta)
 \cong
 Y^+_{\mathrm{Hall}}(X_m;\theta')
\]
must state the permitted changes of stability chamber, grading lattice,
and toric variables. With the toric compact-core grading fixed,
\(s^2=-n\) is an invariant.

\subsection*{Closed statements and stronger identifications}

The closed statements are:
\begin{enumerate}
 \item the toric fan of \(X_n\) has four smooth affine charts
 \(U_i\simeq\C^3\), adjacent overlaps \(\C^2\times\C^*\), and opposite
 overlaps \(\C\times(\C^*)^2\);
 \item \(D_{u_0}\simeq\F_n\) is a compact divisor, so global
 compact-four-cycle-free theorems leave \(X_n\) outside their scope;
 \item \(\CoHA_T(\C^3)=Y^+(\widehat{\mathfrak{gl}}_1)\) on each fixed
 affine chart, while the global local-surface Hall algebra requires the
 \(X_n\) potential and chamber data;
 \item \(\kappa_{\mathrm{ch}}^{\mathrm{loc}}(X_n)=2\), independently
 of \(n\);
 \item the positive Hall half precedes the Drinfeld double, completion,
 derived chiral centre, and any full \(W\)-algebra comparison;
 \item the \(d=3\) \(\Phi\)-output attached to this geometry is
 \(E_1\)-chiral.
\end{enumerate}

A stronger global identification must supply:
\begin{enumerate}
 \item a minimal \((Q_n,W_n)\) or \(A_\infty\) endomorphism model for
 the chosen surface collection, with multiplication relations recording
 \(s^2=-n\);
 \item a PBW or shuffle presentation for \(\CoHA(Q_n,W_n;\theta)\) in
 the compact-four-cycle regime;
 \item a comparison between the refined Hall character and the
 two-variable topological-vertex or Vafa--Witten partition function,
 including the chamber and framing conventions;
 \item named double, centre, completion, and evaluation data before any
 full Yangian or \(W\)-algebra statement.
\end{enumerate}
Local Hirzebruch surfaces therefore form a discrete family of
local-surface Hall positive halves with constant
\(\kappa_{\mathrm{ch}}^{\mathrm{loc}}=2\) and \(n\)-dependent toric and
Hall structure.
